\subsection{The actual Hadamard image of two standard theories}

Let $G$ denote the standard integral PD de Rham theory and
$S=\Gamma_{\mathbb Z}(T)$.  To specify the comparison
precisely, take two copies $G_A,G_B$ and define the
weight-diagonal subalgebra of their tensor over $\mathbb Z$ by
\[
 H_n^{w,q}=\bigoplus_{q_1+q_2=q}
          (G_A)_n^{w,q_1}\otimes_{\mathbb Z}
          (G_B)_n^{w,q_2}.
 \tag{S09.129}
\]
Products, differentials and divided powers are the restrictions
of the ordinary tensor PD algebra.  Equal input weights in
the two factors remain equal under all these operations.
Write $T_A,T_B$ for the coefficient generators of its factors.
The coefficient map $S\to H$ sends
\[
 T^{[w]}\longmapsto
             w!\,T_A^{[w]}\otimes T_B^{[w]}.
 \tag{S09.130}
\]
It is the PD map specified by $T\mapsto T_A\otimes T_B$.

\begin{proposition}[A-S09-26: Proved; the standard Hadamard image and its based closure]
\label{AS0926}
The actual strict maps over the coefficient map (S09.130)
\[
 \begin{split}
 f:G_A\longrightarrow H,&\qquad
             f(a)=a\otimes T_B^w\quad(a\text{ of weight }w),\\
 g:G_B\longrightarrow H,&\qquad
             g(b)=T_A^w\otimes b\quad(b\text{ of weight }w)
 \end{split}                                                     \tag{S09.131}
\]
induce an injective strict map
\[
                    h:G_A\otimes_S G_B\longrightarrow H.
 \tag{S09.132}
\]
Its image $I$ is the PD subalgebra generated by the two
images in (S09.131).  Thus (S09.132) gives an actual
isomorphism $G_A\otimes_S G_B\xrightarrow{\cong}I$.
The entire iterated based PD saturation of $I$ inside
$I\otimes\mathbb Q$, and even inside the larger
$H\otimes\mathbb Q$, leaves $I$ unchanged, with any
weight cutoff.  The vertex-integrality condition uses
the specified coefficient image (S09.130).

For every prime $p$ and every $e\geq1$, the normalized
form-zero row of $I$ has a nonzero class of order $p$
in simplicial homology degree one and weight $p^e$.
Consequently this image, and its displayed based closure,
fail eventual horizontal exactness.
\end{proposition}

\begin{proof}
The maps (S09.131) preserve products because ordinary
powers of $T_A,T_B$ add their exponents.  They preserve
the differential because $dT_A=dT_B=0$.
For each defined divided power of an even homogeneous
element, the tensor PD identity gives
\[
 \gamma_r(a\otimes T_B^w)
       =\gamma_r(a)\otimes T_B^{rw},
 \tag{S09.133}
\]
and the same identity holds for the other factor.
This is the value of $f$ on $\gamma_r(a)$, since its
weight is $rw$.  The addition rule verifies all sums.
The augmentation and simplicial operators are preserved
by the displayed formulas, since coefficient powers
are constant simplicial elements.
Both maps on $T^{[w]}$ equal (S09.130), using
$T^w=w!T^{[w]}$.  They therefore induce the actual
strict map (S09.132) by the tensor universal property.

We prove its injectivity in every simplex and every
bidegree. Split the marked coefficient line in the
standard simplex module.  The resulting integral
description is
\[
 G_n=S\otimes_{\mathbb Z}
       \Gamma_{\mathbb Z}(x_1,\ldots,x_n)
       \otimes_{\mathbb Z}\Lambda_{\mathbb Z}
                              (dx_1,\ldots,dx_n),
 \tag{S09.134}
\]
with all displayed generators of weight one and with
$dx_i=d(x_i)$.  Thus the source of (S09.132) is free
abelian on the corresponding divided-power/exterior
monomials.  It suffices to prove injectivity after
rationalization.

Over $\mathbb Q$, use ordinary polynomial monomials.
For multi-indices $J,L$ of distinct exterior generators,
write $|J|,|L|$ for their sizes.  A source monomial is
\[
              T^a x_A^I dx_A^J x_B^K dx_B^L.
 \tag{S09.135}
\]
Its image is, with the fixed tensor/exterior order,
\[
 T_A^{a+|K|+|L|}\,T_B^{a+|I|+|J|}
                x_A^I dx_A^J x_B^K dx_B^L.
 \tag{S09.136}
\]
The indices $I,J,K,L$ and either coefficient exponent
recover $a$.  Distinct source monomials therefore have
distinct target monomials and are linearly independent.
This proves rational injectivity and hence integral
injectivity of (S09.132).  The tensor algebra is
generated, as a PD algebra, by its two factors, so its
image is exactly the stated $I$.  All maps were already
strict and simplicial; the componentwise inverse onto
the image is consequently strict and simplicial too.

The complex in (S09.134) has an integral coefficient
contraction: for each variable,
$d\gamma_m(x)=\gamma_{m-1}(x)dx$, and the contraction
sends $\gamma_{m-1}(x)dx$ to $\gamma_m(x)$.
Tensoring these finitely many split complexes gives
the coefficient comparison in every weight and simplex.
The same proof, with the two sets of variables, applies
to $G_A\otimes_SG_B$ and hence to $I$.
In particular $H^0(I_n)=S$, $H^1(I_n)=0$, and
$I_0=S$ under (S09.130).
The same contractions after rationalization, tensored
over $\mathbb Q$ at each fixed weight in (S09.129),
give coefficient $H^0$ and vanishing positive vertical
cohomology for $H_{\mathbb Q}$ as well.  Its coefficient
comparison is (S09.130), which is an isomorphism over
$\mathbb Q$ in every weight.

Let $a\in(H_{\mathbb Q})_n^0$ be a rational zero-form
allowed at a based saturation step starting with $I$:
thus $da\in I_n^1$
and the vertex-zero value of $a$ belongs to $S$.
The closed one-form $da$ has an integral primitive
$b\in I_n^0$ by $H^1(I_n)=0$.
The rational coefficient comparison for $H$ says
$a-b\in S_{\mathbb Q}$.  Its vertex value is integral,
because those of both $a$ and $b$ are integral.
Therefore $a-b\in S$ and $a\in I_n^0$.
No new zero-form is added.  The algebra is already
closed under divided powers, so the next stage equals
$I$.  Induction proves the assertion for every stage
and their union, with any cutoff.

We independently compute the horizontal defect.
The degree-zero row of (S09.134) is levelwise free
over $S$; its normalized terms are direct summands
and hence projective over $S$.
The simplicial module of barycentric coordinates has
normalized complex the identity disk
$\mathbb Z[1]\to\mathbb Z[0]$.
Its actual simplicial contraction is preserved by
each positive divided-power functor: applying a functor
to the homotopy maps preserves all composition identities,
and these reduced functors preserve the zero map.
Thus $G^0$ has horizontal homology $\mathbb Z$ in
weight and simplicial degree zero and zero elsewhere.
It is an actual projective $S$-resolution of
$\mathbb Z=S/I_1$, where $I_1=S^{>0}$.

The Eilenberg--Zilber chain equivalence over $S$
therefore computes
\[
 H_1N(G^0\otimes_SG^0)
    =\operatorname{Tor}_1^S(\mathbb Z,\mathbb Z)
    =I_1/I_1^2.
 \tag{S09.137}
\]
The last equality follows by tensoring the exact
sequence $0\to I_1\to S\to\mathbb Z\to0$ and using
the free middle term.  In weight $m>1$ this quotient
is the actual cyclic group
\[
 \mathbb Z\Big/
       \gcd_{0<j<m}\binom mj\,\mathbb Z.
 \tag{S09.138}
\]
For $m=p^e$, every internal binomial coefficient is
divisible by $p$, from $(1+X)^{p^e}=1+X^{p^e}$
over $\mathbb F_p$.  The coefficient at $j=p^{e-1}$
equals $p\binom{p^e-1}{p^{e-1}-1}$, whose second
factor is prime to $p$: all digits of $p^e-1$ are
$p-1$.  No other prime $\ell$ divides all coefficients,
since the base-$\ell$ expansion of $(1+X)^{p^e}$ has
an internal nonzero term when $p^e$ is not a power
of $\ell$.  Thus the gcd is exactly $p$.
Equations (S09.132) and (S09.137) give the claimed
class in $I$, in every such weight.
These weights are unbounded, proving the failure of
eventual horizontal exactness for the specified image
and its unchanged based closure.
\end{proof}

The proposition computes this image and this closure
by actual strict comparisons.  Adjoining further
horizontal primitives is a separate unresolved step
in the main construction.
